\documentclass[11pt,a4paper]{article}
\usepackage[T1]{fontenc}
\usepackage[utf8]{inputenc}
\usepackage[margin=2.2cm]{geometry}
\usepackage{amsmath,amssymb}
\usepackage{booktabs}
\usepackage{listings}
\usepackage{xcolor}
\usepackage[hidelinks]{hyperref}

\lstdefinelanguage{Rust}{
  keywords={let,mut,fn,for,in,if,else,match,move,return,loop,while,break,continue,struct,enum,
            impl,pub,use,true,false,Some,None,Ok,Err,as,ref,usize,bool,i32,u8,Vec,Option},
  sensitive=true,
  comment=[l]{//},
  morestring=[b]",
}
\lstset{
  language=Rust,
  basicstyle=\ttfamily\footnotesize,
  keywordstyle=\color{blue!60!black},
  commentstyle=\color{gray},
  stringstyle=\color{green!40!black},
  columns=fullflexible,
  keepspaces=true,
  breaklines=true,
  frame=single,
  framerule=0.3pt,
  xleftmargin=4pt,
  showstringspaces=false,
}
\newcommand{\code}[1]{\texttt{#1}}

\title{Benchmarks from the related work, written against the Must API\\
\large Specification and implementation pairs for TraceForge, for the implementer}
\author{The author, for the owner}
\date{2026-10-26}

\begin{document}
\maketitle

\section*{What this file is}

The related work (\code{plan/RELATED-AND-BENCHMARKS-TODO.md}, verified section) evaluates on a
small set of protocols: two-phase commit (Disel, Aneris, Trillium), a load balancer (Aneris),
key-value stores (IronKV, Verdi, Grove, Chapar, Memcached in B\'il\'y et al.), leader election
(Igloo), Paxos (IronRSL, Trillium, MoDist's MPS), Raft (Verdi), a lock server (Verdi's
repository). Our five families already cover two-phase commit, a key-value store, a lock, a load
balancer and a leader election. This file gives six more pairs, the realistic protocols the
related work verifies and we do not yet run, each as a specification and an implementation
written in the style of the existing fixtures in \code{conformance/grid.rs}: threads spawned
with \code{named}, typed channels from \code{typed\_chan}, \code{send\_msg},
\code{recv\_msg\_block} (blocking), \code{recv\_msg} (non-blocking, \code{Option}),
\code{nondet()} for a boolean choice and \code{(0..n).nondet()} for a choice in a range,
thread-addressed \code{send\_msg(tid, m)} and \code{recv\_msg\_block::<T>()} where a thread has
no channel. The communication model is the run's (\code{asyn}, \code{p2p} or \code{cd}); each
pair says which it needs. Every program is loop-free after unrolling its bounded parameters,
has static threads, and declares its visible threads by name; the two programs of a pair have
the same visible threads.

The programs are written for the lead to transcribe into fixtures, not to compile as they
stand: they omit the \code{prog(move || \{ \ldots \})} wrapper, the message enums'
derive lines, and the plumbing functions that the existing A1--A4 fixtures already have
(\code{reply\_channels}, \code{spawn\_clients}). Where a fixture needs a choice the Must API does
not have, the text says so. For each pair: the source in the related work, what the two sides
do, where they differ, the visible threads, the parameters and the smallest sizes, the faults to
seed and the verdict expected for each, and what makes the pair worth running.

\paragraph{Four rules every pair follows} (from the review of the first version of this file,
Opus, 2026-10-26; the first version broke three of them).
\begin{enumerate}
  \item \textbf{A non-blocking receive may always read $\bot$}, whether or not a message is
        pending. So an invisible thread that serves clients from a bounded \code{recv\_msg}
        loop has graphs in which it serves nothing, and the clients it serves block. Either the
        specification can stop at any point too (an \code{if nondet() \{ return; \}} in every
        iteration of its server), or the implementation's servers use blocking receives in the
        idiom of A4 (a \code{loop} that ends blocked when the clients are done, the invisible
        status being unobserved). The pairs below use blocking servers wherever a client's
        progress depends on them.
  \item \textbf{One specification thread orders everything it serves.} Two operations of
        different clients served by one invisible thread are causally ordered in every
        specification graph, through that thread's program order. An implementation that
        serves them through independent threads has graphs with no such order, and those are
        reported (\code{prop:refstrict}). So a specification has one server thread per unit the
        implementation may serve independently: one per key for a sharded store, one per
        replica for a replicated store, one oracle for a decision every client observes.
  \item \textbf{Visible threads send and receive only what they send and receive on the other
        side.} A routing update, a retry or a control message received by a visible client is a
        visible event with no counterpart; it goes to an invisible proxy.
  \item \textbf{A fault must be reachable by the script.} Each fault names the schedule and the
        operations that expose it at the smallest size, and a fault that the script cannot
        expose is a conforming control, said so.
\end{enumerate}
\code{typed\_chan} is a FIFO channel and \code{apps\_config} runs \code{ConsType::FIFO}; a pair
that needs reordering (B1's acceptors, B4's updates) builds its channels with a \code{NoOrder}
builder, which the lead adds to the fixture plumbing.

\tableofcontents

\section{B1 --- Single-decree Paxos against a consensus oracle}

\paragraph{Source.} IronFleet's IronRSL (SOSP 2015), Trillium (POPL 2024, single-decree Paxos
refining Lamport's TLA$^+$ model), MoDist's MPS (NSDI 2009), B\'il\'y et al.\ (OOPSLA 2025, \S5.3).
Design reviewed by Opus and Codex (\code{plan/RELATED-OOPSLA25.md} \S4).

\paragraph{What the pair checks.} The implementation runs one instance of single-decree Paxos:
two proposers with ballots 1 and 2, three acceptors, each proposer carrying the value of one
client. The specification is a consensus oracle: it may decide any proposed value, or nothing,
and tells every learner the same decision. The visible threads are the clients, which send
their proposals, and the learners, which receive the decision. Proposers and acceptors are
invisible. Agreement is what the specification imposes: every learner that receives a decision
receives the same value, and only a proposed one.

\paragraph{Messages.}
\begin{lstlisting}
enum ToProposer { Propose(i32) }                       // client -> proposer
enum ToAcceptor { Prepare(u8), Accept(u8, i32) }        // proposer -> acceptor, ballot first
enum ToProposerFromAcc { Promise(u8, Option<(u8, i32)>), Accepted(u8) }
enum ToLearner  { Decide(i32) }
\end{lstlisting}

\paragraph{Specification: the oracle.}
\begin{lstlisting}
// visible: c0, c1, l0, l1       invisible: oracle
let (tx_o, rx_o) = typed_chan::<ToProposer>();             // both clients send here
let (tx_l, rx_l) = learner_channels(2);                    // one channel per learner
named("oracle", move || {
    let mut vals: Vec<i32> = Vec::new();
    for _ in 0..2 {                                        // at most one proposal per client
        if let Some(ToProposer::Propose(v)) = rx_o.recv_msg() { vals.push(v); }
    }
    if !vals.is_empty() && crate::nondet() {               // decide, or decide nothing
        let v = vals[(0..vals.len()).nondet()];
        for tx in &tx_l { tx.send_msg(ToLearner::Decide(v)); }
    }
});
for i in 0..2 { let tx = tx_o.clone();
    named(&format!("c{i}"), move || tx.send_msg(ToProposer::Propose(10 + i as i32))); }
for (i, rx) in rx_l.into_iter().enumerate() {
    named(&format!("l{i}"), move || { let _d: ToLearner = rx.recv_msg_block(); }); }
\end{lstlisting}
The oracle receives each proposal non-blocking and may leave it unread. That is deliberate:
every proposal the oracle consumes becomes a causal predecessor of the decision, and the
implementation's decision depends causally only on the proposals that reached the winning
proposer, so an oracle that consumed every proposal would impose an order the implementation
does not have and produce false alarms (\code{prop:refstrict}). With non-blocking receives the
oracle has a graph for every subset. A learner blocks when nothing is decided, on both sides.

\paragraph{Implementation: Paxos.}
\begin{lstlisting}
// visible: c0, c1, l0, l1       invisible: p1, p2, a0, a1, a2
let (tx_p, rx_p) = proposer_channels(2);     // client i -> proposer i+1 (ballot i+1)
let (tx_a, rx_a) = acceptor_channels(3);
let (tx_pa, rx_pa) = proposer_reply_channels(2);   // acceptor -> proposer, by ballot
let (tx_l, rx_l) = learner_channels(2);
for b in 1..=2u8 {
    let rx_prop = rx_p.remove(0); let rx_rep = rx_pa.remove(0);   // one receiver per proposer
    let tx_a = tx_a.clone(); let tx_l = tx_l.clone();
    named(&format!("p{b}"), move || {
        let ToProposer::Propose(mine) = rx_prop.recv_msg_block();
        for tx in &tx_a { tx.send_msg(ToAcceptor::Prepare(b)); }
        let mut promises = 0; let mut best: Option<(u8, i32)> = None;
        for _ in 0..3 {                                   // collect what arrives, bounded
            if let Some(ToProposerFromAcc::Promise(bb, acc)) = rx_rep.recv_msg() {
                if bb == b { promises += 1;
                    if let Some((ab, av)) = acc {
                        if best.map_or(true, |(x, _)| ab > x) { best = Some((ab, av)); } } } } }
        if promises < 2 { return; }                      // no quorum: give up
        let v = best.map_or(mine, |(_, av)| av);          // the Paxos rule
        for tx in &tx_a { tx.send_msg(ToAcceptor::Accept(b, v)); }
        let mut accepted = 0;
        for _ in 0..3 {
            if let Some(ToProposerFromAcc::Accepted(bb)) = rx_rep.recv_msg() {
                if bb == b { accepted += 1; } } }
        if accepted >= 2 { for tx in &tx_l { tx.send_msg(ToLearner::Decide(v)); } }
    });
}
for a in 0..3 {
    let rx = rx_a[a]; let tx_pa = tx_pa.clone();
    named(&format!("a{a}"), move || {
        let mut promised: u8 = 0; let mut accepted: Option<(u8, i32)> = None;
        for _ in 0..4 {                                   // at most 2 prepares + 2 accepts
            match rx.recv_msg() {
                Some(ToAcceptor::Prepare(b)) if b > promised => {
                    promised = b; tx_pa[b as usize - 1].send_msg(ToProposerFromAcc::Promise(b, accepted)); }
                Some(ToAcceptor::Accept(b, v)) if b >= promised => {
                    promised = b; accepted = Some((b, v));
                    tx_pa[b as usize - 1].send_msg(ToProposerFromAcc::Accepted(b)); }
                _ => {}                                   // stale ballot or nothing pending
            } }
    });
}
for i in 0..2 { let tx = tx_p[i].clone();
    named(&format!("c{i}"), move || tx.send_msg(ToProposer::Propose(10 + i as i32))); }
for (i, rx) in rx_l.into_iter().enumerate() {
    named(&format!("l{i}"), move || { let _d: ToLearner = rx.recv_msg_block(); }); }
\end{lstlisting}
Two proposers may both reach a decision; Paxos guarantees they decide the same value, so each
learner receives one \code{Decide} and reads the first (the second stays pending, unobserved).
If the learners should see exactly one message, make the proposers send to a single invisible
\code{announcer} that forwards once; the oracle is unchanged.

\paragraph{Differences.} The specification has no ballots, no quorums and no acceptors; its
only nondeterminism is which proposal, if any, is decided. The implementation's nondeterminism is
the message order among proposers and acceptors, which decides whether a proposer reaches a
quorum and whether the second proposer adopts the first's value.

\paragraph{Parameters.} Proposers 2, acceptors 3, clients 2, learners 2, ballots \{1, 2\}; the
acceptors' loop bound 4 (two prepares and two accepts) and the proposers' collection bounds 3
(one per acceptor) are what unrolls the protocol, and a late promise eating one accept-phase
iteration still leaves a quorum. Communication: the acceptor and proposer channels \code{NoOrder}
(a Paxos network reorders), the client and learner channels FIFO. Growth: a third proposer with ballot 3, or a third client, or three
learners.

\paragraph{Faults.}
\begin{enumerate}
  \item \code{ignore-best}: the proposer uses \code{mine} even when \code{best} is \code{Some}
        (drop the Paxos rule). Expected: a run in which ballot 1 decides $v_1$ and ballot 2
        decides $v_2 \neq v_1$; two learners observe different values; uncovered, a signature
        mismatch against every oracle graph.
  \item \code{accept-below-promise}: the acceptor accepts \code{Accept(b, v)} with
        \code{b < promised}. Expected: with promises for ballot 1 from $a_0, a_1$ and their
        accepts delayed, ballot 2 chooses $v_2$ through $a_1, a_2$, then $a_0, a_1$ accept the
        delayed ballot-1 $v_1$; disagreement; uncovered.
  \item \code{silent-learner} (status): the winning proposer sends to one learner only.
        Expected: a learner blocked where the oracle's learners are both done; uncovered.
\end{enumerate}

Reviewed (Opus, 2026-10-26): the rule, the bounds and all three faults check out; no schedule
reports the conforming pair, the decided value's client reaching every learner directly or
through adoption, and the oracle consuming only that client's proposal.

\paragraph{Why it is worth running.} The shared protocol of the closest related work; the
first family where the implementation's invisible nondeterminism (message order among five
invisible threads) is large while the specification's is tiny (a choice among at most two
values), the opposite of the load balancer.

\section{B2 --- Raft log replication, bounded, against an atomic log}

\paragraph{Source.} Verdi (PLDI 2015), whose headline is a verified Raft; Grove's replicated
log abstractions; IronRSL's replicated state machine. The specification is the atomic log that
state-machine replication refines.

\paragraph{What the pair checks.} Three replicas run one term of Raft after a bounded election:
$r_0$ and $r_1$ are candidates, $r_2$ votes for the first request it sees, so the leader is
whichever candidate's request reaches $r_2$ first (there is no split at this size; two
candidates with one voter decide by message order). Clients send commands through an invisible
per-client front-end that fans the command out to every replica, so the client itself sends once;
only the leader appends, replicates with \code{AppendEntries}, commits on a majority of acks and
replies to the client with the committed index; a visible reader asks, through its own front-end,
for the committed log. The specification is one atomic log that may stop serving at any point.

\paragraph{Messages.}
\begin{lstlisting}
enum ToReplica { RequestVote(u8, usize), Vote(u8, bool), Cmd(usize /*client*/, i32),
                 AppendEntries(u8, usize /*leader*/, usize /*index*/, i32), Ack(u8, usize), Read }
enum ToClient  { Committed(usize) }
enum ToReader  { Log(Vec<i32>) }
\end{lstlisting}

\paragraph{Specification: the atomic log.}
\begin{lstlisting}
// visible: c0, c1, reader       invisible: log
let (tx_log, rx_log) = typed_chan::<ToReplica>();
let (tx_c, rx_c) = client_channels(2);
let (tx_r, rx_r) = typed_chan::<ToReader>();
named("log", move || {
    let mut log: Vec<i32> = Vec::new();
    for _ in 0..3 {                                 // two commands and one read, any order
        if crate::nondet() { return; }              // may stop here: the implementation may too (rule 1)
        match rx_log.recv_msg_block() {
            ToReplica::Cmd(c, v) => { log.push(v); tx_c[c].send_msg(ToClient::Committed(log.len())); }
            ToReplica::Read => tx_r.send_msg(ToReader::Log(log.clone())),
            _ => {}
        } }
});
for i in 0..2 { let tx = tx_log.clone(); let rx = rx_c[i];
    named(&format!("c{i}"), move || { tx.send_msg(ToReplica::Cmd(i, 100 + i as i32));
                                       let _r: ToClient = rx.recv_msg_block(); }); }
named("reader", move || { tx_log.send_msg(ToReplica::Read); let _l: ToReader = rx_r.recv_msg_block(); });
\end{lstlisting}
The stop in every iteration gives the specification a graph for every subset of clients left
blocked, which the implementation reaches when a replica's non-blocking receive reads $\bot$
(rule 1). The read may observe zero, one or two commands.

\paragraph{Implementation: Raft, one term.}
\begin{lstlisting}
// visible: c0, c1, reader       invisible: f0, f1, fr (front-ends), r0, r1, r2
let (tx_rep, rx_rep) = replica_channels(3);
let (tx_f, rx_f) = frontend_channels(3);               // client i -> f_i, reader -> fr
let (tx_c, rx_c) = client_channels(2);
let (tx_r, rx_r) = typed_chan::<ToReader>();
for i in 0..3 { let rx = rx_f[i]; let tx_rep = tx_rep.clone();
    named(&format!("f{i}"), move || { let m: ToReplica = rx.recv_msg_block();
        for o in 0..3 { tx_rep[o].send_msg(m.clone()); } }); }          // fan-out (rule 3)
for me in 0..3 {
    let rx = rx_rep[me]; let tx_rep = tx_rep.clone(); let tx_c = tx_c.clone(); let tx_r = tx_r.clone();
    named(&format!("r{me}"), move || {
        let term: u8 = 1; let candidate = me < 2;
        let mut voted = false; let mut votes = if candidate { 1 } else { 0 };
        let mut leader: Option<usize> = None;
        let mut log: Vec<i32> = Vec::new(); let mut client_of: Vec<usize> = Vec::new();
        let mut acks: Vec<usize> = Vec::new(); let mut committed = 0usize;
        let mut seen: [bool; 2] = [false, false];       // one append per client (the fan-out)
        let mut waiting: Vec<(usize, i32)> = Vec::new(); // commands received before a leader is known
        let mut read_waiting = false;
        if candidate { for o in 0..3 { if o != me { tx_rep[o].send_msg(ToReplica::RequestVote(term, me)); } } }
        for _ in 0..BOUND {
            match rx.recv_msg() {
                Some(ToReplica::RequestVote(t, cand)) => {
                    let grant = t == term && !voted && !candidate; if grant { voted = true; }
                    tx_rep[cand].send_msg(ToReplica::Vote(t, grant)); }
                Some(ToReplica::Vote(t, true)) if t == term && candidate && leader.is_none() => {
                    votes += 1; if votes >= 2 { leader = Some(me); } }
                Some(ToReplica::Cmd(c, v)) => waiting.push((c, v)),
                Some(ToReplica::AppendEntries(t, l, idx, v)) if t == term => {
                    leader = Some(l);
                    if log.len() + 1 == idx { log.push(v); }
                    tx_rep[l].send_msg(ToReplica::Ack(t, idx)); }
                Some(ToReplica::Ack(t, idx)) if t == term && leader == Some(me) => {
                    acks[idx - 1] += 1;
                    while committed < log.len() && acks[committed] >= 2 {
                        committed += 1; tx_c[client_of[committed - 1]].send_msg(ToClient::Committed(committed)); } }
                Some(ToReplica::Read) => read_waiting = true,
                _ => {}
            }
            if leader == Some(me) {                         // the leader drains what it holds
                for (c, v) in waiting.drain(..) {
                    if seen[c] { continue; } seen[c] = true;
                    log.push(v); client_of.push(c); acks.push(1);
                    let idx = log.len();
                    for o in 0..3 { if o != me { tx_rep[o].send_msg(ToReplica::AppendEntries(term, me, idx, v)); } } }
                if read_waiting { read_waiting = false; tx_r.send_msg(ToReader::Log(log[..committed].to_vec())); }
            } else if let Some(l) = leader {                // a follower forwards what it holds
                for (c, v) in waiting.drain(..) { tx_rep[l].send_msg(ToReplica::Cmd(c, v)); }
                read_waiting = false;                        // the leader also received the Read
            }
        } });
}
for i in 0..2 { let tx = tx_f[i].clone(); let rx = rx_c[i];
    named(&format!("c{i}"), move || { tx.send_msg(ToReplica::Cmd(i, 100 + i as i32));
                                       let _r: ToClient = rx.recv_msg_block(); }); }
named("reader", move || { tx_f[2].send_msg(ToReplica::Read); let _l: ToReader = rx_r.recv_msg_block(); });
\end{lstlisting}
Commands and reads received before a leader is known are buffered locally, not re-sent, so
the loop bound \code{BOUND} (about 10 at this size, derived by the tester) counts only real
messages. The leader deduplicates the three copies of each command by client id; a follower's
copy is forwarded and deduplicated the same way.

\paragraph{Differences.} The specification has one thread and no election; its choices are the
service order and where it stops. The implementation decides the leader by message order,
replicates, and commits on acks. Reads at the leader return the committed prefix, which the
specification's read is at the time of service.

\paragraph{Parameters.} Replicas 3, clients 2, one command each, one reader, one term.
Communication \code{p2p} (per-link FIFO, which Raft assumes). Growth: two commands per client.

\paragraph{Faults.}
\begin{enumerate}
  \item \code{commit-without-quorum}: the leader replies \code{Committed} on its own append.
        Not observable in one term without failures; a conforming control.
  \item \code{out-of-order-commit}: the leader commits index 2 when its acks arrive before index
        1's. Observable by status only: $c_1$ receives \code{Committed(2)} while $c_0$ blocks
        (index 1's acks read as $\bot$); the specification replying 2 to $c_1$ has already
        replied to $c_0$. Uncovered.
  \item \code{follower-serves-read}: a follower answers \code{Read} from its own log. Observable
        twice: the reader gets \code{[100]} while $c_0$ is still blocked (uncommitted), or the
        losing candidate answers \code{[]} after both commits, where the specification's order
        from the read's send to $c_0$'s receive is missing on the implementation side. Uncovered.
\end{enumerate}

\paragraph{Why it is worth running.} Verdi's protocol; the first family with an election and
replication in one implementation. Reviewed: the first version of this pair was reported in
every graph (clients sending three copies, replicas stopping on $\bot$ with no stop in the
specification); this version follows rules 1 and 3.

\section{B3 --- Sharded key-value store with key migration against per-key stores}

\paragraph{Source.} IronFleet's IronKV (SOSP 2015), a sharded store whose shards move keys
between hosts; Grove's reconfigurable store (SOSP 2023).

\paragraph{What the pair checks.} Two shards hold the two keys; each client sends its requests
through an invisible proxy that holds the client's routing table; an invisible controller moves
key 0 from shard 0 to shard 1 during the run and tells the proxies. A request that reaches a
shard that no longer owns the key is forwarded to the owner; a request that reaches the new
owner before the transfer waits there. The specification is one store thread per key (rule 2:
one thread for both keys would order every pair of operations, which two shards never do).
The property is that the migration is invisible.

\paragraph{Messages and script.} \code{Req::Put(c, k, v)}, \code{Req::Get(c, k)},
\code{Rep::Ack}, \code{Rep::Val(Option<i32>)}; the client script of A2 (\code{a2\_op}: operation
$j$ of client $i$ is a \code{Put} iff $i + j$ is even, on key $(\lfloor j/2 \rfloor + \lfloor i/2
\rfloor) \bmod 2$), $Q = 3$ so that key 0, the migrating key, carries both a \code{Put} and a
\code{Get} from different clients.

\paragraph{Specification: one store per key.}
\begin{lstlisting}
// visible: c0, c1       invisible: store0, store1, x0, x1 (the proxies, unchanged on both sides)
let (tx_st, rx_st) = store_channels(2);           // one channel per key
let (tx_c, rx_c) = client_channels(2);
let (tx_x, rx_x) = proxy_channels(2);
for k in 0..2 { let rx = rx_st[k]; let tx_c = tx_c.clone();
    named(&format!("store{k}"), move || { let mut val: Option<i32> = None;
        loop { match rx.recv_msg_block() {                 // blocking: ends blocked when the clients are done
            Req::Put(c, _, v) => { val = Some(v); tx_c[c].send_msg(Rep::Ack); }
            Req::Get(c, _) => tx_c[c].send_msg(Rep::Val(val)),
        } } }); }
for i in 0..2 { let rx = rx_x[i]; let tx_st = tx_st.clone();
    named(&format!("x{i}"), move || { for _ in 0..Q { let r: Req = rx.recv_msg_block();
        let k = match &r { Req::Put(_, k, _) | Req::Get(_, k) => *k as usize };
        tx_st[k].send_msg(r); } }); }
// clients: a2's script, each sending its Q requests to its proxy tx_x[i] and reading Q replies on rx_c[i]
\end{lstlisting}

\paragraph{Implementation: two shards, proxies with routing tables, a controller.}
\begin{lstlisting}
// visible: c0, c1       invisible: s0, s1, x0, x1, ctl
enum Msg { Req(Req), Move(u8, usize), Transfer(u8, Option<i32>) }
enum ToProxy { Req(Req), Owner(u8, usize) }
let (tx_s, rx_s) = shard_channels(2);              // Msg
let (tx_x, rx_x) = proxy_channels(2);              // ToProxy; the client sends ToProxy::Req
let (tx_c, rx_c) = client_channels(2);
named("ctl", move || { if crate::nondet() {        // migrate key 0 at some point, or never
    tx_s[0].send_msg(Msg::Move(0, 1)); for i in 0..2 { tx_x[i].send_msg(ToProxy::Owner(0, 1)); } } });
for me in 0..2 { let rx = rx_s[me]; let tx_s = tx_s.clone(); let tx_c = tx_c.clone();
    named(&format!("s{me}"), move || {
        let mut map: [Option<i32>; 2] = [None, None];
        let mut owns: [bool; 2] = [me == 0, me == 1];  // s0 owns key 0, s1 owns key 1
        let mut parked: Vec<Req> = Vec::new();         // requests for a key not yet transferred in
        loop { match rx.recv_msg_block() {
            Msg::Req(r) => { let k = key_of(&r) as usize;
                if owns[k] { serve(&mut map, &r, &tx_c); }
                else if me == 1 && k == 0 { parked.push(r); }      // the new owner, before Transfer
                else { tx_s[1 - me].send_msg(Msg::Req(r)); } }     // forward to the owner
            Msg::Move(k, to) => { owns[k as usize] = false;
                tx_s[to].send_msg(Msg::Transfer(k, map[k as usize])); }
            Msg::Transfer(k, v) => { map[k as usize] = v; owns[k as usize] = true;
                for r in parked.drain(..) { serve(&mut map, &r, &tx_c); } }
        } } }); }
for i in 0..2 { let rx = rx_x[i]; let tx_s = tx_s.clone();
    named(&format!("x{i}"), move || { let mut owner: [usize; 2] = [0, 1];
        let mut served = 0;
        while served < Q { match rx.recv_msg_block() {
            ToProxy::Owner(k, to) => owner[k as usize] = to,
            ToProxy::Req(r) => { served += 1; tx_s[owner[key_of(&r) as usize]].send_msg(Msg::Req(r)); }
        } } }); }
// clients: a2's script through tx_x[i], replies on rx_c[i]; serve() applies a Put or Get and replies
\end{lstlisting}
The proxy's \code{while} is bounded by the number of requests, with at most one \code{Owner}
message interleaved, so it unrolls. Under per-link FIFO the \code{Transfer} on the link
$s_0 \to s_1$ precedes every request $s_0$ forwards after \code{Move}, so a forwarded request
never bounces; a request that reaches $s_1$ through a proxy that already learnt the new owner
may still precede the \code{Transfer} and is parked.

\paragraph{Differences.} The specification has one store per key and no migration; the
implementation has two shards, a moving key, stale routing tables, forwarded and parked
requests. Both order only same-key operations, through the owner or through the transfer.

\paragraph{Parameters.} Shards 2, keys 2, clients 2, $Q = 3$, one migration. Communication
\code{p2p}. Growth: three clients, or the migration of both keys.

\paragraph{Faults} (all on key 0, the migrating one, with $Q = 3$ so that $c_0$ puts and $c_1$
gets on it).
\begin{enumerate}
  \item \code{stale-after-move}: $s_0$ keeps serving \code{Get} for key 0 from its copy after
        \code{Move}. Schedule: $c_0$'s \code{Put} is served at $s_1$ after the transfer, then
        $c_1$'s \code{Get} reaches $s_0$ through a stale proxy and is answered from the old copy:
        \code{None} or an older value. Uncovered (value).
  \item \code{lost-transfer}: $s_1$ takes ownership on \code{Move} without waiting for
        \code{Transfer} and serves \code{None}. Uncovered (value) after a \code{Put} at $s_0$.
  \item \code{double-apply}: a forwarded \code{Put} is applied and acked by both shards. The
        client receives two acks; the second is read as the reply to its next request, a value
        mismatch there. Uncovered.
\end{enumerate}

\paragraph{Why it is worth running.} IronKV's protocol, and a key-value pair where the
implementation's extra invisible structure changes during the run. Reviewed: the first
version was reported in every graph (a visible client receiving routing updates; one store
thread for both keys; bounded non-blocking shards) and its faults were unreachable (key 1
migrated, the script never put on it); this version follows rules 1 to 4.

\section{B4 --- Causally consistent replicated key-value store against per-replica causal servers}

\paragraph{Source.} Chapar (POPL 2016), certified causally consistent key-value stores; the
Aneris follow-up on a causally consistent store (POPL 2022).

\paragraph{What the pair checks.} Two replicas, each with one client; a client puts and gets at
its own replica; replicas exchange updates over a reordering network and apply a remote update
only after its dependencies. The specification has one server per replica too (rule 2: a single
oracle would order the two clients' operations), each serving its own client in program order
and receiving the other's updates, but applying a remote update at a moment it chooses, subject
to causality. So the specification's extra freedom is exactly the delay a network may add.

\paragraph{Script.} Client $c_0$: \code{Put(x, 1)}, \code{Put(y, 2)}, \code{Get(y)}. Client
$c_1$: \code{Get(y)}, \code{Get(x)}, \code{Put(x, 3)}. Keys $x = 0$, $y = 1$. A read of $y = 2$
at $r_1$ followed by a read of $x = $ \code{None} is the causal violation the faults produce.

\paragraph{Messages.}
\begin{lstlisting}
enum Op  { Put(u8, i32), Get(u8) }
enum Rep { Ack, Val(Option<i32>) }
struct Update { key: u8, val: i32, origin: usize, deps: [usize; 2] }
enum Msg { Op(Op), Upd(Update) }
\end{lstlisting}

\paragraph{Specification: two causal servers.}
\begin{lstlisting}
// visible: c0, c1       invisible: v0, v1
let (tx_v, rx_v) = server_channels(2);     // Msg; the client's ops and the other server's updates (NoOrder)
let (tx_c, rx_c) = client_channels(2);
for me in 0..2 { let rx = rx_v[me]; let tx_v = tx_v.clone(); let tx_c = tx_c.clone();
    named(&format!("v{me}"), move || {
        let mut view: [Option<i32>; 2] = [None, None]; let mut clock: [usize; 2] = [0, 0];
        let mut pending: Vec<Update> = Vec::new();
        loop {
            // before serving, apply any ready pending update, or not, freely (the network's delay)
            let mut i = 0;
            while i < pending.len() { let u = pending[i].clone();
                let ready = u.deps[u.origin] == clock[u.origin] + 1 && u.deps[1 - u.origin] <= clock[1 - u.origin];
                if ready && crate::nondet() { view[u.key as usize] = Some(u.val); clock[u.origin] += 1; pending.remove(i); }
                else { i += 1; } }
            match rx.recv_msg_block() {
                Msg::Op(Op::Put(k, v)) => { clock[me] += 1; view[k as usize] = Some(v);
                    tx_v[1 - me].send_msg(Msg::Upd(Update { key: k, val: v, origin: me, deps: clock }));
                    tx_c[me].send_msg(Rep::Ack); }
                Msg::Op(Op::Get(k)) => tx_c[me].send_msg(Rep::Val(view[k as usize])),
                Msg::Upd(u) => pending.push(u),
            } } }); }
for i in 0..2 { /* the clients, each sending its three ops to tx_v[i] as Msg::Op and reading three replies */ }
\end{lstlisting}

\paragraph{Implementation: two replicas with dependency tracking.}
\begin{lstlisting}
// visible: c0, c1       invisible: r0, r1
// identical to the servers above with one difference: a pending update whose dependencies are
// met is applied at once (no nondet), the network's reordering being the only source of delay.
\end{lstlisting}

\paragraph{Differences.} The two sides are structurally the same two threads; the
specification may delay a ready update, the implementation applies it at once. Since the
update channel reorders, the implementation's graphs are the specification's graphs in which
the delay is realised by the network, and the specification has more.

\paragraph{Parameters.} Replicas 2, clients 2, three operations each, keys 2. The update
channels \code{NoOrder}; the client channels FIFO. Growth: three replicas with three-entry clocks.

\paragraph{Faults.}
\begin{enumerate}
  \item \code{apply-eagerly}: the replica applies a buffered update without checking its
        dependencies. Schedule: $c_0$ puts $x$ then $y$; the update for $y$ overtakes the one
        for $x$ on the reordering channel; $r_1$ applies $y$ at once; $c_1$ reads $y = 2$ then
        $x = $ \code{None}. The specification can never return that pair. Uncovered (value).
  \item \code{lost-update-with-clock}: the replica advances its clock for a ready update but
        does not write the value. Schedule: as above with the updates in order; $c_1$ reads
        $y = $ \code{None} after $r_1$ has advanced past it, and a later dependent update applies
        on top of the missing one, so the final read differs from every specification graph.
        Uncovered (value). The variant that drops the update without advancing the clock is
        not observable, since the specification may delay forever; it is a conforming control.
\end{enumerate}
The fault \code{wrong-origin-clock} of the first version is dropped: with two replicas it only
delays updates, which the specification allows.

\paragraph{Why it is worth running.} Chapar's property as a specification program; the pair
where the specification is the implementation plus a choice, which makes the index's
containment test do the work. Reviewed: the first version had one oracle (reported in the
common schedule) and a script that never read a written key; this version follows rules 2 and 4.
Estimate revised from 2 to 3 lead days.

\section{B5 --- A Memcached-style server with connections against one map}

\paragraph{Source.} B\'il\'y et al.\ (OOPSLA 2025, \S5.1, Memcached v1--v3 with Get, Set,
Delete over connections, one thread per connection, a store behind a lock). Extends our A2.

\paragraph{What the pair checks.} Each client holds one connection and issues a sequence of
\code{Get}, \code{Set} and \code{Delete} on two keys. The implementation has one worker per
connection; the workers share no memory and reach the store through a storage thread by message
passing, which is where the server's internal concurrency lives. The specification is one map.

\paragraph{Messages.}
\begin{lstlisting}
enum Cmd { Get(u8), Set(u8, i32), Delete(u8) }
enum Res { Val(Option<i32>), Stored, Deleted, NotFound }
\end{lstlisting}

\paragraph{Specification: one map} (as \code{a2\_spec}, with the \code{Delete} arm:
\code{Deleted} if the key was present, \code{NotFound} otherwise; both clients send
\code{(i, Cmd)} to one channel, served in arrival order, the reply to client $i$). Visible: the
clients. Invisible: \code{store}. One store thread is right here (rule 2): the implementation's
single storage thread orders every pair of commands too.

\paragraph{Implementation: a worker per connection and a storage thread.}
\begin{lstlisting}
// visible: c0, c1       invisible: w0, w1, storage
let (tx_w, rx_w) = worker_channels(2);      // client i -> worker i (its connection)
let (tx_st, rx_st) = typed_chan::<(usize, Cmd)>();   // workers -> storage
let (tx_wr, rx_wr) = worker_reply_channels(2);        // storage -> worker i
let (tx_c, rx_c) = client_channels(2);
named("storage", move || {
    let mut map: [Option<i32>; 2] = [None, None];
    for _ in 0..2 * Q {
        let (w, cmd) = rx_st.recv_msg_block();
        let res = match cmd {
            Cmd::Get(k) => Res::Val(map[k as usize]),
            Cmd::Set(k, v) => { map[k as usize] = Some(v); Res::Stored }
            Cmd::Delete(k) => if map[k as usize].take().is_some() { Res::Deleted } else { Res::NotFound },
        };
        tx_wr[w].send_msg(res);
    }
});
for w in 0..2 {
    let rx = rx_w[w]; let tx_st = tx_st.clone(); let rx_wr = rx_wr[w]; let tx_c = tx_c.clone();
    named(&format!("w{w}"), move || { for _ in 0..Q {
        let (c, cmd): (usize, Cmd) = rx.recv_msg_block();   // the same (i, Cmd) the spec receives
        tx_st.send_msg((c, cmd));
        let res: Res = rx_wr.recv_msg_block();
        tx_c[w].send_msg(res);
    } });
}
for i in 0..2 { let tx = tx_w[i].clone(); let rx = rx_c[i];
    named(&format!("c{i}"), move || {                       // both connections on key 0: Set, Get, Delete
        for cmd in [Cmd::Set(0, 10 + i as i32), Cmd::Get(0), Cmd::Delete(0)] {
            tx.send_msg((i, cmd)); let _r: Res = rx.recv_msg_block(); } }); }
\end{lstlisting}

\paragraph{Differences.} The specification serves commands in one arrival order; the
implementation interleaves two workers in front of one storage thread, so the service order is
the storage thread's arrival order, which the workers' relays reshuffle. With a store behind a
single thread the pair conforms; the realistic faults are in the workers.

\paragraph{Parameters.} Clients 2, $Q = 3$ (\code{Set}, \code{Get}, \code{Delete}, both on key
0, so that the connections interact). Communication \code{p2p} (TCP connections). Growth: three
clients, or a second key with a fourth command.

\paragraph{Faults.}
\begin{enumerate}
  \item \code{cached-get}: a worker answers a \code{Get} on key $k$ from the value its own
        connection last \code{Set} on $k$, without asking the storage. Schedule: $c_0$ sets
        $10$, $c_1$ sets $11$, $c_0$ gets and is answered $10$ from the worker's cache while the
        storage holds $11$; the specification serving $c_0$'s \code{Get} after $c_1$'s \code{Set}
        returns $11$, and serving it before has the order $c_0$'s get before $c_1$'s set, absent
        from the implementation graph. Uncovered.
  \item \code{delete-ack-early}: the worker replies \code{Deleted} without asking the storage.
        Schedule: $c_0$ sets, $c_0$ deletes (\code{Deleted}), $c_1$ deletes and is also answered
        \code{Deleted} by its worker, where the storage would say \code{NotFound}; no
        specification order gives two \code{Deleted} after one \code{Set}. Uncovered.
\end{enumerate}
The fault \code{reorder-connection} of the first version is dropped: the clients are
closed-loop, so a worker never holds two commands.

\paragraph{Why it is worth running.} The case study of the closest related work, written in
our language; against their 118 to 181 lines of Rust and 286 to 377 lines of proof annotations,
our pair is two programs of this size and no annotations.

\section{B6 --- Ring leader election against an oracle naming the maximum}

\paragraph{Source.} Igloo (OOPSLA 2020) verifies a leader election protocol among its three case
studies; the ring election of Chang and Roberts is the standard one. We have a leader election
already (A5, within two-phase commit); this one is the classic ring and stands alone.

\paragraph{What the pair checks.} $n$ nodes in a ring, each with a unique identifier; each node
sends its identifier to its successor; a node forwards an identifier larger than its own, drops
a smaller one, and the node that receives its own identifier is the leader and sends
\code{Leader(id)} around the ring once. Every node has a visible observer that receives the
announcement. The specification is an oracle that sends \code{Leader(max id)} to every observer.

\paragraph{Specification.}
\begin{lstlisting}
// visible: o0..o(n-1)       invisible: oracle
let (tx_o, rx_o) = observer_channels(n);
named("oracle", move || { let m = IDS.iter().max().unwrap(); for tx in &tx_o { tx.send_msg(Leader(*m)); } });
for (i, rx) in rx_o.into_iter().enumerate() { named(&format!("o{i}"), move || { let _l: Leader = rx.recv_msg_block(); }); }
\end{lstlisting}

\paragraph{Implementation: Chang--Roberts.}
\begin{lstlisting}
// visible: o0..o(n-1)       invisible: n0..n(n-1)
let (tx_n, rx_n) = node_channels(n);
let (tx_o, rx_o) = observer_channels(n);
for me in 0..n {
    let rx = rx_n[me]; let succ = tx_n[(me + 1) % n].clone(); let tx_obs = tx_o[me].clone();
    named(&format!("n{me}"), move || {
        let my = IDS[me];
        succ.send_msg(Ring::Elect(my));
        for _ in 0..n {                                   // at most n election messages pass a node
            match rx.recv_msg_block() {
                Ring::Elect(id) if id > my => succ.send_msg(Ring::Elect(id)),
                Ring::Elect(id) if id == my => { succ.send_msg(Ring::Leader(my)); tx_obs.send_msg(Leader(my)); return; }
                Ring::Elect(_) => {}                       // smaller: drop
                Ring::Leader(id) => { tx_obs.send_msg(Leader(id)); if id != my { succ.send_msg(Ring::Leader(id)); } return; }
            } }
    });
}
for (i, rx) in rx_o.into_iter().enumerate() { named(&format!("o{i}"), move || { let _l: Leader = rx.recv_msg_block(); }); }
\end{lstlisting}
With \code{IDS} a fixed permutation, e.g.\ $[2, 0, 3, 1]$. The loop bound $n$ covers the
longest path an identifier travels; the leader's own \code{Leader} message returns to it and
the \code{return} ends the thread.

\paragraph{Differences.} The specification knows the maximum; the implementation discovers it
by forwarding, with the message interleavings of $n$ nodes. The observable is one value per
observer and the statuses.

\paragraph{Parameters.} $n \in \{3, 4, 5\}$. Communication \code{p2p} (a ring link is FIFO).
Growth: $n = 6$, and the variant with two identifiers equal (a tie the protocol must not have).

\paragraph{Faults.}
\begin{enumerate}
  \item \code{forward-smaller}: the node holding the maximum forwards a smaller identifier
        instead of dropping it. Expected: the second-largest identifier returns to its owner,
        which announces itself; observers see a wrong leader; uncovered (value). Placed at any
        other node the fault is a conforming control.
  \item \code{drop-announcement}: a node other than the leader's predecessor does not forward
        \code{Leader}. Expected: the observers beyond it block; uncovered (status). At the
        leader's predecessor the announcement has already gone round, a conforming control.
  \item \code{announce-twice}: the leader forwards its own \code{Leader} message again on
        return. Expected: an observer receives two announcements where the oracle sends one;
        the second is pending and unobserved, so this is a conforming control at this size.
\end{enumerate}

\paragraph{Why it is worth running.} Igloo's protocol; the cheapest pair in this file (one
afternoon), and one where the specification has a single graph. Reviewed: correct under
per-link FIFO, the bound $n$ tight and sufficient.

\section{Summary and order of implementation}

\begin{center}
\begin{tabular}{llllll}
\toprule
pair & protocol & source & smallest size & model & lead days \\
\midrule
B1 & single-decree Paxos & IronFleet, Trillium, MoDist, OOPSLA'25 & 2 proposers, 3 acceptors, 2 clients, 2 learners & NoOrder & 2 \\
B2 & Raft replication, one term & Verdi & 3 replicas, 2 clients, 1 reader & p2p & 4 \\
B3 & sharded store with migration & IronKV, Grove & 2 shards, 2 clients, $Q=3$ & p2p & 3 \\
B4 & causally consistent store & Chapar & 2 replicas, 2 clients, 3 ops each & NoOrder & 3 \\
B5 & Memcached-style server & OOPSLA'25 & 2 clients, $Q=3$ & p2p & 1 \\
B6 & ring leader election & Igloo & $n=4$ & p2p & 0.5 \\
\bottomrule
\end{tabular}
\end{center}

Order by value for the paper: B1, B6, B5, B3, B2, B4. B1 is the shared protocol of the closest
related work and reviewed sound; B6 is cheap and reviewed sound; B5 is their case study in our
language; B3 and B2 are the realistic protocols of the deductive line, each three to four days;
B4 last, since only two of its faults are observable and its specification is nearly its
implementation.
Each pair goes to the lead as one item: the tester derives the expected verdicts and uncovered
sets at the smallest size by hand before any run, the conforming pair and every fault are run
under both engines at the smallest size and one size up, and the digest records, as for A1--A5,
the graph counts of both sides, the wall times, the peaks and the first-report times. Nothing
here runs before the final campaign plan names it.
\end{document}
